% STATUS: DRAFT
\chapter{Large $N$ and emergent type III$_1$}\label{ch:large_N}
\skippable

\begin{physmotiv}
In AdS/CFT, a boundary theory at large $N$ describes bulk gravity. The boundary algebra of a single-sided region is type I in finite $N$ regularizations but, in the large-$N$ limit around a thermal state (the eternal black hole), the algebra of single-trace operators becomes a type III$_1$ algebra: an emergent type III$_1$ factor with an emergent time, the Schwarzschild time of the exterior. This is the black hole counterpart of the type III$_1$ static patch algebra, and the starting point for the crossed-product treatment of the black hole. Papers: Leutheusser and Liu on emergent times and subalgebra--subregion duality; Witten, and Chandrasekaran--Penington--Witten on large $N$ algebras.
\end{physmotiv}

\section{The two-sided black hole and the thermofield double}

The eternal AdS black hole is dual to the thermofield double of two copies of the CFT, $\ket{\rm TFD}=Z^{-1/2}\sum\ee^{-\beta E_n/2}\ket n_L\ket n_R$ (\cref{ch:tomita}). At finite $N$ the right CFT algebra is $\Bnd{\HH_R}$, type I, and the TFD is a purification of the thermal state with $\Delta=\ee^{-\beta(H_R-H_L)}$. In the strict large-$N$ limit, define the algebra $\mathcal M_R$ generated by (bounded functions of) single-trace operators smeared in time, in the GNS representation of the thermal state. Because the thermal state has an infinite number of low-lying single-particle excitations of the bulk fields in the $N\to\infty$ limit, the modular spectrum becomes continuous, and the algebra is expected to be type III$_1$ (Leutheusser--Liu), the bulk free fields in the exterior being a generalized free field theory at $N=\infty$ with a continuous spectrum for $\Delta$, with the boundary time translation as outer modular flow. In this limit, boundary time $\leftrightarrow$ modular time up to $\beta$, as in \cref{ch:rs_bw}.

\section{What this tells us}
\begin{itemize}
\item \textbf{Emergent horizon.} The type III$_1$ factor encodes the bulk exterior; its commutant is the bulk exterior of the other side, the algebra of $L$ (Haag duality at large $N$).
\item \textbf{Why type III.} The bulk exterior has horizon-ward modes at all energy scales. Type III$_1$ is the expected type of the exterior bulk algebra for the same reason as in \cref{ch:typeIII_local}.
\item \textbf{Subalgebra--subregion duality.} Different boundary subalgebras correspond to different bulk regions; this is the algebraic form of entanglement-wedge reconstruction (\cref{ch:ew_qec}).
\end{itemize}

\begin{whatwrong}
Type III$_1$ appears only in the large-$N$ limit: at finite $N$ the boundary algebra is type I. The statement that this is type III$_1$ relies on properties of generalized free field theory at $N=\infty$ and is argued, not proved for any interacting CFT. The entropy of the boundary theory at finite $N$ is finite; the large-$N$ type III$_1$ limit loses it in the same way as the continuum limit of a lattice.
\end{whatwrong}

\section*{Exercises}
\begin{exercise}
For a single bulk free scalar in the BTZ-like exterior with quasinormal-mode density, argue that the thermal state of the single-trace algebra has $\Delta$ with spectrum $\ee^{-\beta\omega}$, $\omega\in\R$, hence continuous.
\end{exercise}
\begin{exercise}
Explain why the commutant of the right single-trace algebra in the TFD GNS space is the left single-trace algebra.
\end{exercise}
